\documentclass[11pt,a4paper]{article}

\usepackage[utf8]{inputenc}
\usepackage[T1]{fontenc}
\usepackage{amsmath, amssymb, amsthm}
\usepackage{graphicx}
\usepackage{geometry}
\usepackage{hyperref}
\usepackage{enumitem}
\usepackage{booktabs}
\usepackage{fancyhdr}

\geometry{margin=2.5cm}
\hypersetup{colorlinks=true,linkcolor=blue,citecolor=blue,urlcolor=blue}

% ==== الفوتر الشامل - بعد ازالة v0 ====
\pagestyle{fancy}
\fancyhf{}
\fancyfoot[C]{\tiny
Abdulla Al-Jabri — \href{mailto:jabri62018@gmail.com}{jabri62018@gmail.com} |
V1 \href{https://doi.org/10.5281/zenodo.20100622}{10.5281/zenodo.20100622} |
V2 \href{https://doi.org/10.5281/zenodo.23241084}{10.5281/zenodo.23241084} |
OSF \href{https://doi.org/10.17605/OSF.IO/YC9MB}{10.17605/OSF.IO/YC9MB} |
\href{https://github.com/Jabri-web/Jabri-com}{GitHub} |
\href{https://github.com/Jabri-web/Jabri-com/wiki}{Wiki} |
\href{https://explore.openaire.eu/search/result?pid=10.5281/zenodo.23241084}{OpenAIRE}
}
\fancyhead[C]{}
\renewcommand{\headrulewidth}{0pt}
\renewcommand{\footrulewidth}{0.4pt}

\title{
    On the Impossibility of Reproducing Quantum Correlations \\
    with a Real-Valued Mother Function $Z(x)$: \\
    A Constraint on Future Extensions
}
\author{
    Abdulla Mohammed Nasser Al-Jabri\thanks{Independent Researcher, Sana'a, Yemen.
    ORCID: \href{https://orcid.org/0009-0003-3319-3822}{0009-0003-3319-3822}.
    Email: \href{mailto:jabri62018@gmail.com}{jabri62018@gmail.com}.
    Project: \href{https://osf.io/yc9mb/}{osf.io/yc9mb} —
    Code: \href{https://github.com/Jabri-web/Jabri-com}{github.com/Jabri-web/Jabri-com} —
    Wiki: \href{https://github.com/Jabri-web/Jabri-com/wiki}{github.com/Jabri-web/Jabri-com/wiki}}
}
\date{October 2026}

\begin{document}
\maketitle
\thispagestyle{fancy}

\begin{abstract}
\noindent
This short note reports a negative but informative result concerning the
so-called \emph{Mother Function} $Z(x) = x^5 \ln(x) \sin(2\pi/x) e^{-x/x_p}$,
proposed as part of the Sundus Unified Theory ($Z + C + A = 1$).
We construct a local hidden-variable model directly from $Z(x)$ and test
it against the CHSH inequality. The model yields $S = 1.82$, respecting
Bell's local bound $S \leq 2$. We prove that this is not a numerical
accident: any model in which the measurement outcome depends only on a
shared real variable $x$ and a real phase $\sin(2\pi/x)$ is necessarily
local, and therefore cannot reproduce the quantum value $S = 2\sqrt{2}$
observed in Aspect-type experiments. This establishes a rigorous
constraint: any future extension of $Z(x)$ aiming to capture quantum
entanglement must possess (i) a complex-valued structure that enters the
probability density and (ii) a non-factorizable bilinear form
$Z(x_1, x_2) \neq Z(x_1) Z(x_2)$.

The numerical value $S_Z = 1.82$ is not a discovery; it is the expected
consequence of Bell's theorem applied to a local hidden-variable model.
\end{abstract}

\vspace{1em}
\noindent\textbf{Keywords:} Bell's theorem, CHSH inequality, quantum entanglement,
local hidden variables, Mother Function, Sundus Theory.

\section{Introduction}
The \emph{Sundus Unified Theory} proposes $Z + C + A = 1$, in which $Z(x)$
plays the role of structural generator. We ask whether $Z(x)$ can encode
quantum entanglement without non-local signalling. This note answers in
the negative for the real-valued form and shows it is a necessary consequence
of Bell's theorem.

\section{Preliminaries}
\begin{equation}
    Z(x) = x^5 \ln(x) \sin\!\left(\frac{2\pi}{x}\right) e^{-x/x_p},\quad x>0
\end{equation}
$x$ is dimensionless (in units of $l_p$). $x_p=1.22$ was used for this numerical
evaluation only, as a dimensionless scale parameter; no physical unit is claimed.
\begin{equation}
    \phi(x):=2\pi/x,\quad m_a(x)=\mathrm{sign}[\sin(\phi(x)+a)],\quad
    p(x)=\frac{|Z(x)|^2}{\int |Z|^2 dx}
\end{equation}

\section{Construction}
\begin{equation}
    E_Z(a,b)=\int_0^\infty p(x) m_a(x) m_b(x) dx
\end{equation}
This is a local hidden-variable model: outcome depends only on shared $x$ and local setting.

\begin{table}[h]
\centering
\begin{tabular}{cccc}
\toprule
$\theta$ (deg) & $E_{\mathrm{QM}}=-\cos\theta$ & $E_Z$ (this work) & $|diff|$ \\
\midrule
0 & $-1.0000$ & $+1.0000$ & $2.0000$ \\
22.5 & $-0.9239$ & $+1.0000$ & $1.9239$ \\
45 & $-0.7071$ & $+0.9989$ & $1.7060$ \\
67.5 & $-0.3827$ & $+0.9882$ & $1.3709$ \\
90 & $0.0000$ & $+0.9104$ & $0.9104$ \\
112.5 & $+0.3827$ & $+0.5055$ & $0.1228$ \\
135 & $+0.7071$ & $-0.5674$ & $1.2745$ \\
157.5 & $+0.9239$ & $-0.9997$ & $1.9236$ \\
180 & $+1.0000$ & $-1.0000$ & $2.0000$ \\
\bottomrule
\end{tabular}
\caption{Non-circular test: $x\in[0.05,30]$, $2\times10^5$ points, $S_Z=1.82$}
\end{table}

\section{Main Result}
$S=E(a,b)-E(a,b')+E(a',b)+E(a',b')$, $a=0,a'=\pi/2,b=\pi/4,b'=-\pi/4$.
$S_Z=1.82<2$. Any model $E(a,b)=\int p(x)m_a(x)m_b(x)dx$ satisfies $|S|\le2$ (Bell-CHSH).

The numerical value $S_Z = 1.82$ is not a discovery; it is the expected
consequence of Bell's theorem applied to a local hidden-variable model.

\section{Why it fails}
(i) No antisymmetry: singlet needs $E=-\cos(a-b)$. Real $Z(x)$ symmetric.
$Z_{new}=x^5\ln(x)e^{i\phi}e^{-x/x_p}$ still gives $|Z_{new}|^2=|Z|^2$ → local.
(ii) Factorizability: $Z(x)$ is single-variable, entanglement needs $Z(x_1,x_2)\neq Z(x_1)Z(x_2)$.

\section{Necessary conditions}
(i) Complex structure entering $p(x)$, not only amplitude.
(ii) Non-factorizable: $Z(x_1,x_2)=Z(x_1)Z(x_2)e^{i\Phi(x_1,x_2)}$, $\Phi\neq\Phi_1+\Phi_2$.

\section{Conclusion}
Real-valued $Z(x)$ cannot reproduce quantum correlations. $S_Z=1.82$ is Bell's bound, not failure. This constrains future extensions.

\begin{thebibliography}{9}
\bibitem{Bell1964} J.~S.~Bell, Physics \textbf{1}, 195 (1964).
\bibitem{CHSH1969} J.~F.~Clauser, M.~A.~Horne, A.~Shimony, and R.~A.~Holt, Phys. Rev. Lett. \textbf{23}, 880 (1969).
\bibitem{Aspect1982} A.~Aspect et al., Phys. Rev. Lett. \textbf{49}, 1804 (1982).
\bibitem{Nobel2022} Nobel Prize in Physics 2022.
\end{thebibliography}
\end{document}